\section{Conclusion}\label{sec:conclusion}
Evaluating a forecast for a fixed-volume execution problem requires attention to the prices available across the execution horizon and to the costs of reallocating the order. A terminal forecast does not identify those temporal contrasts unless additional restrictions connect it to the full price path. The matched-endpoint diagnostic makes that omission explicit while holding terminal-return scores exactly constant. Its economic accounting then asks a separate question: whether the realised benefit of the schedule change pays for the action it induces.

The Bitcoin--Ether illustration shows why these questions should be separated. The learned profile has positive average gross timing alignment under the primary specification, but curvature cost is larger and full exposure produces a negative net objective gain. Historical attenuation changes the balance, yet the reported plug-in interval includes zero and the lower-percentile rule chooses baseline abstention throughout. The sign of the pooled learned gain changes when April is excluded in an ex post descriptive check, showing that the negative mean is concentrated rather than stable across months. Neither policy establishes a generally superior execution procedure. The negative result is informative because the decomposition identifies where the within-model loss arises, rather than attributing it to endpoint accuracy alone.

The evidence supports a practical reporting recommendation: pair the forecast's stated statistical target with its induced schedule, timing payoff, action cost and baseline-relative gain. Those quantities should be distinguished even when a method is not profitable, and their economic scale should be reported alongside uncertainty. Here the primary loss is approximately 0.44 parts per million of reference notional, so statistical detectability cannot establish commercial importance.

A stronger operational assessment would require explicit order sizes, quotes or depth, estimated execution charges, and genuinely prospective evaluation across additional regimes. Competing learned temporal profiles would test whether the diagnostic changes model selection among plausible alternatives. Dependence-sensitive inference and repeated training would address different uncertainties from the present conditional day bootstrap. These extensions could change the empirical performance conclusions; they are not results already obtained. The current study provides a reproducible, deliberately bounded basis for asking those questions, without presenting familiar execution algebra as a new theory or hypothetical gains as attainable trading savings.
